\section{Conclusion}
\label{sec:conclusion}

We introduced \benchname, a benchmark of 22 ordering tasks spanning five property categories, and used it to probe which features LLMs can most reliably access from parametric memory.
Our experiments with \gptlarge and \gptsmall reveal a property-accessibility hierarchy: algorithmic operations and memorized sequences ($\tau = 1.0$) are trivially easy, historical chronologies are near-perfect ($\tau \approx 0.97$), approximate magnitudes are strong but imperfect ($\tau \approx 0.93$), and precise numeric facts are hardest ($\tau \approx 0.83$).
The critical distinction is not whether knowledge is well-known or obscure, but whether it is stored as a \emph{sequence} or as \emph{independent numeric values}.

These findings have direct practical implications: LLM-based data pipelines should use external tools for ordering by precise numeric properties, while trusting LLMs for sequential and categorical orderings.
Scientifically, our results suggest that LLMs encode factual knowledge in at least two forms---sequential chains that support reliable ordering, and approximate numeric representations that support coarse but not fine-grained comparison.

Future work should scale to longer lists, test additional model families, investigate whether chain-of-thought reasoning can bridge the gap on hard tasks, and explore whether ordering accuracy correlates with fact frequency in training data.
The \benchname benchmark and property-accessibility hierarchy provide a foundation for this line of inquiry.
